\subsection{معماری و پیاده‌سازی}
پیاده‌سازی این سیستم در فایل \lr{Codes/framework/compression\_engine/huffman\_coder.py} انجام یافته و در رجیستری فشرده‌سازی با نام‌های \lr{huffman} و \lr{huffman\_coder} به ثبت رسیده است. ویژگی‌های معماری موتور شامل موارد زیر است:
\begin{itemize}
    \item \textbf{گره و درخت هافمن (\lr{HuffmanNode} و \lr{HuffmanCoder}):} ماژول تخصصی جهت محاسبه برداری فراوانی نمادها، ساخت سریع درخت پیشوندی در حافظه و بسته‌بندی بیت‌ها در آرایه‌های فشرده \lr{numpy.uint8}.
    \item \textbf{کوانتایش پیوسته وزن‌ها (\lr{Uniform Weight Quantization}):} نگاشت متقارن مقادیر ممیز شناور به کدهای صحیح با تعداد بیت‌های قابل تنظیم ($b \in \{4, 8\}$).
    \item \textbf{تثبیت وزن پایه و حفظ گرادیان \lr{LoRA}:} تابع \lr{\_freeze\_base\_only} تضمین می‌کند که پارامترهای اصلی کوانتایز شده تثبیت گردیده و تنها پارامترهای تنظیمی \lr{LoRA} در طول بازیابی گرادیانی آموزش‌پذیر باقی بمانند.
    \item \textbf{حسابرسی دقیق بایت‌ها (\lr{Exact Byte Accounting}):} محاسبه نرخ فشرده‌سازی فیزیکی بر اساس تعداد بایت‌های واقعی فشرده‌شده و جداول کدهای ذخیره‌شده.
\end{itemize}

\subsection{محدودیت‌ها و تحلیل شکست}
تحلیل عملکردی ماژول نقاط ضعف زیر را نمایان می‌سازد:
\begin{itemize}
    \item \textbf{افت دقت ناشی از کوانتایش تکمیلی در مدل‌های عمیق (\lr{RoBERTa} و \lr{GPT-2}):} کدگذاری هافمن به خودی خود بدون اتلاف است، اما اتلاف دقت از مرحله پیش‌نیاز کوانتایش گسسته ناشی می‌شود. در مدل‌های با لایه‌های متوالی متعدد، تجمع خطای کوانتایش در تمام لایه‌ها موجب افت دقت به سطح پایه تصادفی شده و وضعیت \lr{PASS\_LOW} را ایجاد می‌کند.
    \item \textbf{ناهمگونی تاخیر پردازشی در دیکودینگ ناهمگن (\lr{Variable Bit-Length Decoding}):} به علت طول متغیر کدهای باینری، دیکودینگ موازی در سطح رشته‌های موازی (\lr{SIMD/Warp}) در پردازنده‌های گرافیکی دچار انحراف انشعاب (\lr{Warp Divergence}) شده و نیازمند پیاده‌سازی کرنل‌های اختصاصی سخت‌افزاری است.
\end{itemize}
